% GENERATED FILE. Edit canonical lessons, then run npm run generate:books.
% canonical-source-hash: fnv1a64:501ecea2ef8d7770
% canonical-lessons: LA-C271-populus, LA-C271-civitas, LA-C271-patria, LA-C271-hostis, LA-C271-bellum

\chapter{People, State, Homeland, Enemy, War}
\label{ch:la-people-state-homeland-enemy-war}
\hlchaptermodality{271}

\begin{chapteropening}
\textbf{By the end of this chapter:} \emph{I can say a people, a state, a homeland, an enemy and war.}

Complete the last lesson of chapter 271: I can say a people, a state, a homeland, an enemy and war.
\end{chapteropening}

\section[populus, populī]{populus, populī — a people, a nation}

\label{lesson:LA-C271-populus}



\begin{quote}
Before the new one: say the Latin for an anchor, then the Latin for an oar.
\end{quote}

\subsection*{You'll want to know: populus, populī}
\textbf{populus, populī} — \textquotedblleft{}a people, a nation\textquotedblright{}.

\begin{cousinweb}
\textbf{populus} has no sure origin; EtymDB 2.1 gives Proto-Italic \emph{*poplos}, \textquotedblleft{}army\textquotedblright{}, for its oldest sense, a body of armed men. It may be a loan, perhaps from Etruscan. English \emph{people} comes from it through French.
\end{cousinweb}

\begin{grammarlens}[title={What you've built}]
One more word for this chapter.
\end{grammarlens}

\subsection*{Your turn}
\begin{itemize}
\raggedright
  \item \emph{Say it:} \emph{populus, populī}
  \item \emph{Say it:} \emph{populus, populī}, once more
  \item \emph{Say it:} say \emph{populus, populī}
  \item \emph{Recall:} say the Latin for an anchor, then the Latin for an oar, then say \emph{populus, populī} again
\end{itemize}

\begin{tcolorbox}[breakable,colback=teal!4,colframe=teal!35!black,title={Before you move on}]
What does populus, populī mean? (\textquotedblleft{}A people, a nation\textquotedblright{}.) Say it once more.
\end{tcolorbox}
\section[cīvitās, cīvitātis]{cīvitās, cīvitātis — a state; citizenship}

\label{lesson:LA-C271-civitas}



\begin{quote}
Before the new one: say the Latin for an oar, then the Latin for a people.
\end{quote}

\subsection*{You'll want to know: cīvitās, cīvitātis}
\textbf{cīvitās, cīvitātis} — \textquotedblleft{}a state; citizenship\textquotedblright{}.

\begin{cousinweb}
\textbf{cīvitās} is built on \emph{cīvis}, \textquotedblleft{}a citizen\textquotedblright{} (Lewis \& Short), with \emph{-tās}, an ending for a quality: the body of citizens. English \emph{city} comes from it through French.
\end{cousinweb}

\begin{grammarlens}[title={What you've built}]
One more word for this chapter.
\end{grammarlens}

\subsection*{Your turn}
\begin{itemize}
\raggedright
  \item \emph{Say it:} \emph{cīvitās, cīvitātis}
  \item \emph{Say it:} \emph{cīvitās, cīvitātis}, once more
  \item \emph{Say it:} say \emph{cīvitās, cīvitātis}
  \item \emph{Recall:} say the Latin for an oar, then the Latin for a people, then say \emph{cīvitās, cīvitātis} again
\end{itemize}

\begin{tcolorbox}[breakable,colback=teal!4,colframe=teal!35!black,title={Before you move on}]
What does cīvitās, cīvitātis mean? (\textquotedblleft{}A state; citizenship\textquotedblright{}.) Say it once more.
\end{tcolorbox}
\section[patria, patriae]{patria, patriae — a homeland}

\label{lesson:LA-C271-patria}



\begin{quote}
Before the new one: say the Latin for a people, then the Latin for a state; citizenship.
\end{quote}

\subsection*{You'll want to know: patria, patriae}
\textbf{patria, patriae} — \textquotedblleft{}a homeland\textquotedblright{}.

\begin{cousinweb}
\textbf{patria} is the feminine of \emph{patrius}, \textquotedblleft{}of a father\textquotedblright{} (Lewis \& Short), with \emph{terra}, \textquotedblleft{}land\textquotedblright{}, understood: the land of one's fathers.
\end{cousinweb}

\begin{grammarlens}[title={What you've built}]
One more word for this chapter.
\end{grammarlens}

\subsection*{Your turn}
\begin{itemize}
\raggedright
  \item \emph{Say it:} \emph{patria, patriae}
  \item \emph{Say it:} \emph{patria, patriae}, once more
  \item \emph{Say it:} say \emph{patria, patriae}
  \item \emph{Recall:} say the Latin for a people, then the Latin for a state; citizenship, then say \emph{patria, patriae} again
\end{itemize}

\begin{tcolorbox}[breakable,colback=teal!4,colframe=teal!35!black,title={Before you move on}]
What does patria, patriae mean? (\textquotedblleft{}A homeland\textquotedblright{}.) Say it once more.
\end{tcolorbox}
\section[hostis, hostis]{hostis, hostis — an enemy (of the state)}

\label{lesson:LA-C271-hostis}



\begin{quote}
Before the new one: say the Latin for a state; citizenship, then the Latin for a homeland.
\end{quote}

\subsection*{You'll want to know: hostis, hostis}
\textbf{hostis, hostis} — \textquotedblleft{}an enemy (of the state)\textquotedblright{}.

\begin{cousinweb}
\textbf{hostis} is inherited, usually traced to Proto-Indo-European \emph{*ghostis}, \textquotedblleft{}a stranger, a guest\textquotedblright{} (EtymDB 2.1): the stranger became the enemy in Latin, the guest in Germanic. English \emph{guest} is a cousin.
\end{cousinweb}

\begin{grammarlens}[title={What you've built}]
One more word for this chapter.
\end{grammarlens}

\subsection*{Your turn}
\begin{itemize}
\raggedright
  \item \emph{Say it:} \emph{hostis, hostis}
  \item \emph{Say it:} \emph{hostis, hostis}, once more
  \item \emph{Say it:} say \emph{hostis, hostis}
  \item \emph{Recall:} say the Latin for a state; citizenship, then the Latin for a homeland, then say \emph{hostis, hostis} again
\end{itemize}

\begin{tcolorbox}[breakable,colback=teal!4,colframe=teal!35!black,title={Before you move on}]
What does hostis, hostis mean? (\textquotedblleft{}An enemy (of the state)\textquotedblright{}.) Say it once more.
\end{tcolorbox}
\section[bellum, bellī]{bellum, bellī — war}

\label{lesson:LA-C271-bellum}



\begin{quote}
Before the new one: say the Latin for a homeland, then the Latin for an enemy.
\end{quote}

\subsection*{You'll want to know: bellum, bellī}
\textbf{bellum, bellī} — \textquotedblleft{}war\textquotedblright{}.

\begin{cousinweb}
\textbf{bellum} comes from older \emph{duellum} (Lewis \& Short); Latin turned an initial dw into b, as in \emph{bis}, \textquotedblleft{}two times\textquotedblright{}, from older \emph{duis}. The deeper origin of \emph{duellum} is uncertain.
\end{cousinweb}

\begin{grammarlens}[title={What you've built}]
One more word for this chapter.
\end{grammarlens}

\subsection*{Your turn}
\begin{itemize}
\raggedright
  \item \emph{Say it:} \emph{bellum, bellī}
  \item \emph{Say it:} \emph{bellum, bellī}, once more
  \item \emph{Say it:} say \emph{bellum, bellī}
  \item \emph{Recall:} say the Latin for a homeland, then the Latin for an enemy, then say \emph{bellum, bellī} again
\end{itemize}

\begin{tcolorbox}[breakable,colback=teal!4,colframe=teal!35!black,title={Before you move on}]
What does bellum, bellī mean? (\textquotedblleft{}War\textquotedblright{}.) Say it once more.
\end{tcolorbox}
